%! Carrying Out a Test for the Difference of Two Population Means — Unit 7, Lesson 7.9 — Homework
\documentclass[10pt]{article}
\usepackage{apstats-article}
\usepackage{apstats-boxes}
\usepackage{grading}
\gradingassignment{apstats-u07-l09-hw}
\gradingcourse{advanced_statistics}
\newcommand{\TallMath}[1]{$\displaystyle #1\rule[-1.4em]{0pt}{3.2em}$}

\begin{document}

\pageheader{Unit 7, Lesson 7.9}{Homework}

\vspace{0.1in}

\textit{Complete all four SPDC steps for each problem. Show all computations.}

\begin{enumerate}

\item \textbf{Study Hours.} \quad
A university administrator wants to know whether students who live on campus study more
hours per week on average than commuter students. Random samples of 35 residential students
($\bar x_1 = 18.4$ h, $s_1 = 5.2$ h) and 30 commuter students ($\bar x_2 = 15.9$ h,
$s_2 = 6.1$ h) are selected. Dotplots for both groups are roughly symmetric with no outliers.

\begin{enumerate}[label=\alph*.]
  \item \textbf{State:} $H_0$: \makebox[4cm]{\smash{\gradedblank[id=hw1a1, type=short, points=1, answer={{mu1 - mu2 = 0}}, check=text, display={\boldmath$\mu_1 - \mu_2 = 0$}]{3.8cm}{0.19in}}} \quad $H_a$: \makebox[4cm]{\smash{\gradedblank[id=hw1a2, type=short, points=1, answer={{mu1 - mu2 > 0}}, check=text, display={\boldmath$\mu_1 - \mu_2 > 0$}]{3.8cm}{0.19in}}}
  \boxguard
  \item \textbf{Plan:} Method: \makebox[5cm]{\smash{\gradedblank[id=hw1b1, type=short, points=1, answer={{two-sample t-test for mu1 - mu2}}, check=text, display={\boldmath two-sample $t$-test for $\mu_1 - \mu_2$}]{4.8cm}{0.19in}}}, conditions met (state briefly):\par
        \gradedhide{\gradedhere{\gradedblank[id=hw1b2, type=rubric, points=1, display={\bfseries\boldmath random; CLT}]{\linewidth+6pt}{0.19in}}\writelines{9}}
  \item \textbf{Do:} $SE = \sqrt{\dfrac{s_1^2}{n_1} + \dfrac{s_2^2}{n_2}} \approx$ \makebox[2.5cm]{\smash{\gradedblank[id=hw1c1, type=short, points=1, answer=1.419, check=numeric, tolerance=0.01, display={\boldmath$1.419$}]{2.3cm}{0.19in}}}
        \quad $t \approx$ \makebox[2cm]{\smash{\gradedblank[id=hw1c2, type=short, points=1, answer=1.76, check=numeric, tolerance=0.02, display={\boldmath$1.76$}]{1.8cm}{0.19in}}} \quad Technology: $p \approx$ \makebox[2cm]{\smash{\gradedblank[id=hw1c3, type=short, points=1, answer=0.042, check=numeric, tolerance=0.003, display={\boldmath$0.042$}]{1.8cm}{0.19in}}}, $df \approx$ \makebox[2cm]{\smash{\gradedblank[id=hw1c4, type=short, points=1, answer=57.4, check=numeric, tolerance=0.6, display={\boldmath$57.4$}]{1.8cm}{0.19in}}}
  \boxguard
  \item \textbf{Conclude:} At $\alpha = 0.05$: \gradedhide{\writeline\\\gradedhere{\gradedblank[id=hw1d, type=rubric, points=2, display={\bfseries\boldmath $p \approx 0.042 < 0.05$: reject $H_0$; convincing evidence residential students study more on average.}]{\linewidth+6pt}{\gradedlinesheight{1}}}\writelines{9}}
\end{enumerate}

\vspace{0.05in}

\newpage
\item \textbf{Reading Scores.} \quad
An education researcher wants to know if two different literacy programs produce different
mean reading scores for third graders. Program A: $n_1 = 40$, $\bar x_1 = 72.3$, $s_1 = 9.8$.
Program B: $n_2 = 38$, $\bar x_2 = 68.7$, $s_2 = 11.2$. Students were randomly assigned
to programs; CLT applies.

\begin{enumerate}[label=\alph*.]
  \item \textbf{State:} $H_0$: \makebox[4cm]{\smash{\gradedblank[id=hw2a1, type=short, points=1, answer={{mu1 - mu2 = 0}}, check=text, display={\boldmath$\mu_1 - \mu_2 = 0$}]{3.8cm}{0.19in}}} \quad $H_a$: \makebox[4cm]{\smash{\gradedblank[id=hw2a2, type=short, points=1, answer={{mu1 - mu2 != 0}}, check=text, display={\boldmath$\mu_1 - \mu_2 \ne 0$}]{3.8cm}{0.19in}}}
  \boxguard
  \item \textbf{Plan:} Method: \makebox[5cm]{\smash{\gradedblank[id=hw2b1, type=short, points=1, answer={{two-sample t-test for mu1 - mu2}}, check=text, display={\boldmath two-sample $t$-test for $\mu_1 - \mu_2$}]{4.8cm}{0.19in}}}. Conditions (brief):\par
        \gradedhide{\gradedhere{\gradedblank[id=hw2b2, type=rubric, points=1, display={\bfseries\boldmath random; $n$'s $\ge 30$}]{\linewidth+6pt}{0.19in}}\writelines{9}}
  \item \textbf{Do:} $SE \approx$ \makebox[2.5cm]{\smash{\gradedblank[id=hw2c1, type=short, points=1, answer=2.388, check=numeric, tolerance=0.01, display={\boldmath$2.388$}]{2.3cm}{0.19in}}} \quad $t \approx$ \makebox[2cm]{\smash{\gradedblank[id=hw2c2, type=short, points=1, answer=1.51, check=numeric, tolerance=0.02, display={\boldmath$1.51$}]{1.8cm}{0.19in}}}
        \quad $p$ (two-sided) $\approx$ \makebox[2cm]{\smash{\gradedblank[id=hw2c3, type=short, points=1, answer=0.136, check=numeric, tolerance=0.003, display={\boldmath$0.136$}]{1.8cm}{0.19in}}} \quad $df \approx$ \makebox[2cm]{\smash{\gradedblank[id=hw2c4, type=short, points=1, answer=73.5, check=numeric, tolerance=0.6, display={\boldmath$73.5$}]{1.8cm}{0.19in}}}
  \boxguard
  \item \textbf{Conclude:} \gradedhide{\writeline\\\gradedhere{\gradedblank[id=hw2d, type=rubric, points=2, display={\bfseries\boldmath $p \approx 0.136 > 0.05$: fail to reject $H_0$; no convincing evidence the programs' mean scores differ.}]{\linewidth+6pt}{\gradedlinesheight{1}}}\writelines{9}}
\end{enumerate}

\vspace{0.05in}

\newpage
\item \textbf{Recovery Time.} \quad
A physical therapist claims a new stretching protocol reduces recovery time after a knee
injury. A randomized experiment assigns patients to the new protocol ($n_1 = 18$,
$\bar x_1 = 14.2$ days, $s_1 = 3.6$ days) or the standard protocol ($n_2 = 20$,
$\bar x_2 = 17.5$ days, $s_2 = 4.1$ days). Both distributions are roughly symmetric with
no outliers.

\begin{enumerate}[label=\alph*.]
  \item \textbf{State:} $H_0$: \makebox[4cm]{\smash{\gradedblank[id=hw3a1, type=short, points=1, answer={{mu1 - mu2 = 0}}, check=text, display={\boldmath$\mu_1 - \mu_2 = 0$}]{3.8cm}{0.19in}}} \quad $H_a$: \makebox[4cm]{\smash{\gradedblank[id=hw3a2, type=short, points=1, answer={{mu1 - mu2 < 0}}, check=text, display={\boldmath$\mu_1 - \mu_2 < 0$}]{3.8cm}{0.19in}}}
  \boxguard
  \item \textbf{Plan:} Method: \makebox[5cm]{\smash{\gradedblank[id=hw3b1, type=short, points=1, answer={{two-sample t-test for mu1 - mu2}}, check=text, display={\boldmath two-sample $t$-test for $\mu_1 - \mu_2$}]{4.8cm}{0.19in}}}. Conditions:
        \begin{itemize}
          \boxguard
          \item Random: \hfill\makebox[4.6in][l]{\gradedhere{\gradedblank[id=hw3b2, type=rubric, points=1, display={\bfseries\boldmath randomized experiment: random assignment --- met}]{4.6in}{0.19in}}}\\
                \gradedhide{\writelines{9}}
          \boxguard
          \item Normal ($n < 30$): \hfill\makebox[4.6in][l]{\gradedhere{\gradedblank[id=hw3b3, type=rubric, points=1, display={\bfseries\boldmath both roughly symmetric, no outliers --- met}]{4.6in}{0.19in}}}\\
                \gradedhide{\writelines{9}}
        \end{itemize}
  \item \textbf{Do:} $SE \approx$ \makebox[2.5cm]{\smash{\gradedblank[id=hw3c1, type=short, points=1, answer=1.249, check=numeric, tolerance=0.01, display={\boldmath$1.249$}]{2.3cm}{0.19in}}} \quad $t \approx$ \makebox[2cm]{\smash{\gradedblank[id=hw3c2, type=short, points=1, answer=-2.64, check=numeric, tolerance=0.02, display={\boldmath$-2.64$}]{1.8cm}{0.19in}}}
        \quad Technology: $p \approx$ \makebox[2cm]{\smash{\gradedblank[id=hw3c3, type=short, points=1, answer=0.006, check=numeric, tolerance=0.002, display={\boldmath$0.006$}]{1.8cm}{0.19in}}} \quad $df \approx$ \makebox[2cm]{\smash{\gradedblank[id=hw3c4, type=short, points=1, answer=36.0, check=numeric, tolerance=0.6, display={\boldmath$36.0$}]{1.8cm}{0.19in}}}
  \boxguard
  \item \textbf{Conclude:} At $\alpha = 0.05$: \gradedhide{\writeline\\\gradedhere{\gradedblank[id=hw3d, type=rubric, points=2, display={\bfseries\boldmath $p \approx 0.006 < 0.05$: reject $H_0$; convincing evidence the new protocol lowers mean recovery time.}]{\linewidth+6pt}{\gradedlinesheight{1}}}\writelines{9}}
\end{enumerate}

\vspace{0.05in}

\newpage
\item \textbf{Fuel Efficiency.} \quad
An automotive engineer tests whether a new engine design produces a different mean fuel
efficiency (mpg) than the current design. Current: $n_1 = 25$, $\bar x_1 = 31.4$ mpg,
$s_1 = 2.8$ mpg. New: $n_2 = 25$, $\bar x_2 = 33.1$ mpg, $s_2 = 3.2$ mpg. Random samples;
dotplots are roughly symmetric.

\begin{enumerate}[label=\alph*.]
  \item \textbf{State:} $H_0$: \makebox[4cm]{\smash{\gradedblank[id=hw4a1, type=short, points=1, answer={{mu1 - mu2 = 0}}, check=text, display={\boldmath$\mu_1 - \mu_2 = 0$}]{3.8cm}{0.19in}}} \quad $H_a$: \makebox[4cm]{\smash{\gradedblank[id=hw4a2, type=short, points=1, answer={{mu1 - mu2 != 0}}, check=text, display={\boldmath$\mu_1 - \mu_2 \ne 0$}]{3.8cm}{0.19in}}}
  \boxguard
  \item \textbf{Plan:} Method: \makebox[5cm]{\smash{\gradedblank[id=hw4b1, type=short, points=1, answer={{two-sample t-test for mu1 - mu2}}, check=text, display={\boldmath two-sample $t$-test for $\mu_1 - \mu_2$}]{4.8cm}{0.19in}}}. Conditions (brief):\par
        \gradedhide{\gradedhere{\gradedblank[id=hw4b2, type=rubric, points=1, display={\bfseries\boldmath random; symmetric}]{\linewidth+6pt}{0.19in}}\writelines{9}}
  \item \textbf{Do:} $SE \approx$ \makebox[2.5cm]{\smash{\gradedblank[id=hw4c1, type=short, points=1, answer=0.850, check=numeric, tolerance=0.01, display={\boldmath$0.850$}]{2.3cm}{0.19in}}} \quad $t \approx$ \makebox[2cm]{\smash{\gradedblank[id=hw4c2, type=short, points=1, answer=-2.00, check=numeric, tolerance=0.02, display={\boldmath$-2.00$}]{1.8cm}{0.19in}}}
        \quad $p$ (two-sided) $\approx$ \makebox[2cm]{\smash{\gradedblank[id=hw4c3, type=short, points=1, answer=0.051, check=numeric, tolerance=0.002, display={\boldmath$0.051$}]{1.8cm}{0.19in}}}
  \boxguard
  \item \textbf{Conclude:} \gradedhide{\writeline\\\gradedhere{\gradedblank[id=hw4d, type=rubric, points=2, display={\bfseries\boldmath $p \approx 0.051 > 0.05$: fail to reject $H_0$; no convincing evidence the mean mpg differs.}]{\linewidth+6pt}{\gradedlinesheight{1}}}\writelines{9}}
\end{enumerate}

\end{enumerate}

\begin{reflectionbox}
A student says: ``We got $p = 0.18$, which means the two groups probably have equal means.''
Explain what is wrong with this statement. What does ``fail to reject $H_0$'' actually mean
when comparing two groups?
\writelines{3}
\end{reflectionbox}

\newpage
\begin{extensionbox}
\textbf{Preview of Lesson 7.10 --- Selecting an Inference Procedure.}
For each scenario below, identify the correct inference procedure (one-sample $t$-test,
matched-pairs $t$-test, or two-sample $t$-test) and explain your reasoning.
\begin{enumerate}[label=(\alph*)]
  \item Researchers measure blood pressure for the same 15 patients before and after
        taking a new medication.
  \item Two independent groups of students take different versions of a standardized test.
        Researchers compare mean scores.
  \item A factory samples 40 bolts and tests whether the mean diameter differs from the
        specified 10 mm.
\end{enumerate}
\writelines{4}
\end{extensionbox}

\end{document}
